\subsubsection{Raster- versus punktwolkenbasierte Verfahren zur Einzelbaumsegmentierung}

Bei der Segmentierung von Bäumen auf Basis von LiDAR-Daten wird grundlegend zwischen \textit{raster- und punktwolkenbasierten Verfahren} unterschieden. Im Folgenden werden die beiden Ansätze gegenübergestellt und deren Eignung für den Kontext urbaner Bäume eingeordnet.

Beim \textit{rasterbasierten Segmentierungsverfahren} wird die 3D-Punktwolke zunächst auf ein planares Raster projiziert, um ein Kronenhöhenmodell (Canopy Height Model, \ac{CHM}) oder ein \textit{digitales Oberflächenmodell (\ac{DOM})} zu erzeugen. Auf dieses 2D-Raster können Bildverarbeitungsalgorithmen angewandt werden. Ein typischer Ablauf beginnt mit der Identifizierung von Baumspitzen über Local Maximum Filter. Diese detektierten Spitzen dienen als Startpunkte für anschließendes \textit{Regio Growing}, bei dem benachbarte Pixel nach bestimmten Regeln hinzufügt werden. Eine weitere gängige Methode ist der \textit{Watershed-} bzw. \textit{Marker-controlled Watershed-Algorithmus}, welcher das Fließen von Wasser auf einem invertierten Rastermodell simuliert, um Baumgrenzen zu isolieren. Darauf wird im Folgenden noch genauer eingegangen. Ergänzend Ansätze umfassen \textit{Valley Following}, welcher gezielt nach lokalen Minima sucht, um Lücken und Vertiefungen zwischen einzelnen Baumkronen zu detektieren sowie \textit{Graph-Cut}, der entlang der schwächsten Kante die Kronengrenzen schneidet \cite{fuIndividualTreeSegmentationUAV2024}.

Der wesentliche Vorteil dieses Ansatzes liegt in seiner hohen Recheneffizienz, wodurch er sich für großflächige Analysen eignet. Zudem lässt sich ein CHM mit geringem Aufwand generieren. Demgegenüber steht jedoch ein systematischer Verlust an 3D-Informationen im Rasterisierungsprozess, da primär die obere Kronenschicht abgebildet wird. Infolgedessen bleiben Bäume in Unterschichten unberücksichtigt, was insbesondere in dichten, mehrschichtigen Beständen zu Segmentierungsfehlern führt \cite{munzingerMappingUrbanForest2022}. Zudem können bei der Generierung des Rastermodells Fehlstellen (\textit{Pits}) entstehen, welche die Genauigkeit nachfolgender Algorithmen stark beeinträchtigen und eine vorherige Interpolation bzw. Glättung erfordern \cite{liIndividualTreeSegmentation2023}.

Beim \textit{punktwolkenbasierten Ansatz} werden dagegen direkt die 3D-Daten verarbeitet, entweder als Rohdaten oder in einer diskretisierten Voxelstruktur. Da dieses Verfahren auf der vollständigen räumlichen Geometrie operiert, ermöglicht es auch die Detektion von Bäumen in niedrigeren Bestandschichten \cite{munzingerMappingUrbanForest2022}. Algorithmisch wird dieser Ansatz entweder durch \textit{3D-Regio Growing} umgesetzt, bei dem Punkte basierend auf räumlichen Distanzen und Normalenvektoren zu Segmenten verschmolzen werden oder durch \textit{Clustering-Methoden}. Zu diesen zählt unter anderem \textit{K-Means-Clustering}, ein partitionsbasierter Algorithmus, welcher die Punktwolke in eine festgelegte Anzahl an Clustern aufteilt, indem er die Quadrate der euklidischen Distanz zwischen den Punkten und den jeweiligen Clusterschwerpunkten minimiert. Ein flexiblerer dichtebasierte Ansatz wird durch \textit{Density-Based Spatial Clustering of Applications with Noise (\ac{DBSCAN})} realisiert. Hier werden Punkte zu Clustern zusammengefasst, wenn sie innerhalb eines definierten Suchradius eine Mindestanzahl an nachbarpunkten aufweisen \cite{fuIndividualTreeSegmentationUAV2024}.

Trotz der höheren Informationsdichte führt dieser Ansatz nicht zwangsläufig zu besseren Ergebnissen als rasterbasierte Verfahren. Klassische Clustering-Algorithmen erweisen sich oft als ungeeignet für heterogene Bestände mit stark variierenden Baumgrößen und unregelmäßigen Abständen. \textit{Punktbasierte Regio Growing-Verfahren} erfordern meist wissensbasierte, feste Annahmen über morphologische Baumstrukturen \cite{fuIndividualTreeSegmentationUAV2024}. Diese komplexen Parameterabhängigkeiten schränken die Übertragbarkeit der Modelle auf andere Baumbestände stark ein. Darüber hinaus geht die Verarbeitung der 3D-Punkte mit einem erheblichen Rechen- und Speicheraufwand einher \cite{liIndividualTreeSegmentation2023}. Welche Methode geeigneter ist, hängt somit stark vom Untersuchungsgebiet ab. Für diese Umsetzung wurde der rasterbasierter Ansatz gewählt auf dem im Methodikkapitel genauer eingegangen wird.
 
